\documentclass[aps,prl,twocolumn,superscriptaddress]{revtex4-2}
\usepackage{amsmath,amssymb,graphicx,hyperref}

\begin{document}

\title{BCT Appendix App-L209a:\\
Aquifer Detection via BCT Bessel Ratio Mapping of SAR Backscatter}
\author{Michel Robert Cabri\'e}
\email{ZeroFreeParameters@gmail.com}
\affiliation{Independent Artist and Researcher, Victoria, Australia\\
ORCID: 0009-0007-9561-9859}
\date{March 2026}

\begin{abstract}
We derive the expected BCT Bessel ratio signatures for fresh water aquifers in Sentinel-1 C-band SAR data. Water-saturated porous rock produces a distinctive ratio $R_{\mathrm{BCT}} = 0.75 \pm 0.10$ due to the dielectric contrast between wet ($\epsilon_r \approx 25$) and dry ($\epsilon_r \approx 6$) geological media. The ratio is distinguishable from mineralised rock ($R_{\mathrm{BCT}} \approx 1.0$), dry rock ($R_{\mathrm{BCT}} \approx 0.45$), and clay ($R_{\mathrm{BCT}} \approx 0.55$). Depth sensitivity in arid soils extends to 2--5~m at C-band, sufficient for shallow aquifer detection. A processing pipeline using Google Earth Engine is described. Total cost: zero. BCT Prediction \#256. Zero free parameters.
\end{abstract}

\maketitle

\section{Dielectric Contrast}

SAR backscatter intensity is governed by the complex dielectric permittivity of the surface and near-surface materials. For geological media at C-band (5.405~GHz, Sentinel-1):

\begin{table}[h]
\centering
\begin{tabular}{lcc}
\hline
Material & $\epsilon_r$ & $R_{\mathrm{BCT}}$ (predicted) \\
\hline
Dry sand & 3--5 & 0.35--0.45 \\
Dry rock (granite) & 5--8 & 0.45--0.55 \\
Clay (dry) & 5--10 & 0.50--0.60 \\
Wet soil & 15--25 & 0.65--0.80 \\
Water-saturated rock & 20--30 & 0.70--0.85 \\
Mineralised quartz & 4--6 (with metallic inclusions) & 0.90--1.10 \\
Standing water & 80 & $>$1.2 \\
\hline
\end{tabular}
\caption{Dielectric permittivity and predicted BCT ratio for geological materials at C-band.}
\end{table}

The BCT ratio $R_{\mathrm{BCT}}$ is computed as:
\begin{equation}
R_{\mathrm{BCT}} = \frac{\sigma^0_{\mathrm{pixel}}}{\langle \sigma^0 \rangle_{\mathrm{local}}} \cdot \frac{r_{\mathrm{tet}}}{r_{\mathrm{oct}}} = \frac{\sigma^0_{\mathrm{pixel}}}{\langle \sigma^0 \rangle_{\mathrm{local}}} \times 0.5426
\end{equation}

The normalisation by $r_{\mathrm{tet}}/r_{\mathrm{oct}}$ centres the scale such that BCT-resonant materials (mineralised rock) produce $R_{\mathrm{BCT}} \to 1$, while non-resonant materials produce $R_{\mathrm{BCT}} < 1$.

\section{Aquifer Signature}

A fresh water aquifer manifests as a spatially coherent region where $R_{\mathrm{BCT}}$ shifts from the background dry-rock value ($\sim$0.45) to the wet-rock value ($\sim$0.75). The signature characteristics:

\begin{enumerate}
\item \textbf{Spatial coherence:} Aquifers extend over areas $>$1~km$^2$, producing smooth $R_{\mathrm{BCT}}$ anomalies rather than point targets.
\item \textbf{Temporal stability:} Unlike surface moisture (which varies with rainfall), aquifer signatures persist across dry seasons. Multi-temporal stacking enhances detectability.
\item \textbf{Depth gradient:} The $R_{\mathrm{BCT}}$ anomaly increases toward the centre of the aquifer where the water table is shallowest, producing a characteristic `dome' shape in the ratio map.
\item \textbf{Boundary sharpness:} Aquifer edges produce steep $R_{\mathrm{BCT}}$ gradients where saturated and unsaturated zones meet.
\end{enumerate}

\section{Depth Sensitivity}

C-band radar (5.6~cm wavelength) penetrates the surface to a depth dependent on soil moisture and composition:
\begin{equation}
d_{\mathrm{penetration}} \approx \frac{\lambda}{2\pi} \cdot \sqrt{\frac{\epsilon_r}{\epsilon''}} \approx 0.5\text{--}5~\text{m}
\end{equation}
where $\epsilon''$ is the imaginary part of the dielectric constant (loss factor).

In arid environments (Sahel, Australian outback, Middle East), low soil moisture reduces $\epsilon''$, extending penetration depth to 3--5~m. This is sufficient to detect the upper surface of shallow aquifers, which represent the most accessible and valuable water resources in water-stressed regions.

For deeper aquifers, L-band SAR (23.6~cm wavelength, available from ALOS-2 and future NISAR) extends penetration to 10--15~m.

\section{Processing Pipeline}

The complete processing chain uses freely available tools:

\begin{enumerate}
\item \textbf{Data:} Sentinel-1 GRD (Ground Range Detected) products from Copernicus Open Access Hub. Free. Global coverage.
\item \textbf{Platform:} Google Earth Engine. Free for research. Handles petabyte-scale SAR data natively.
\item \textbf{Preprocessing:} Thermal noise removal, orbit correction, radiometric calibration, terrain correction (all standard GEE functions).
\item \textbf{BCT Bessel ratio computation:} Apply Eq.~(2) at each pixel using a local averaging kernel of radius 500~m.
\item \textbf{Multi-temporal stacking:} Average $R_{\mathrm{BCT}}$ across all available dry-season acquisitions (typically 20--50 scenes per year) to suppress speckle and enhance persistent anomalies.
\item \textbf{Thresholding:} Flag pixels with $0.65 < R_{\mathrm{BCT}} < 0.85$ as candidate aquifer zones.
\item \textbf{Spatial filtering:} Retain only spatially coherent anomalies $>$0.5~km$^2$.
\end{enumerate}

Total computational cost in Google Earth Engine: minutes per 10,000~km$^2$ survey area.

\section{Validation Strategy}

Each target class requires ground truth. For aquifers:
\begin{itemize}
\item Compare BCT SAR aquifer predictions with existing hydrogeological survey data (borehole logs, resistivity surveys)
\item Partner with hydrogeologists in water-stressed regions (Sahel, Horn of Africa, Australian outback, Middle East)
\item Drill test boreholes at predicted vs.\ non-predicted locations (controlled validation)
\end{itemize}

Initial validation can use existing borehole databases (e.g., Bureau of Meteorology groundwater database for Australia, BGS borehole records for UK) at zero cost.

\section{Humanitarian Impact}

The United Nations estimates 2.2 billion people lack access to safely managed drinking water~\cite{WHO2023}. The majority live in regions with adequate groundwater resources that are unmapped due to the cost of conventional hydrogeological surveys (\$10,000--\$100,000 per survey area).

BCT SAR aquifer mapping costs zero. The satellite data is free. The processing platform is free. The algorithm is three numbers. A national aquifer survey for any country on Earth can be completed in days, not years.

Under BCT-EPL-1.0, this capability is free for all nations, forever.

\begin{acknowledgments}
Dedicated to the memory of Robert Cabri\'e, Bureau of Meteorology, Antarctica 1966. He understood that the Earth's water belongs to everyone.

And to the 2.2 billion people who need it most. The geometry is yours. It always was.
\end{acknowledgments}

\begin{thebibliography}{6}
\bibitem{WHO2023} World Health Organization, Progress on household drinking water, sanitation and hygiene 2000--2022 (WHO/UNICEF, 2023).
\bibitem{L116} M.~R.~Cabri\'e, BCT Letter 116: The Lost Archive, Zenodo (2026).
\bibitem{L209} M.~R.~Cabri\'e, BCT Letter 209: The Eye of the Lattice, Zenodo (2026).
\bibitem{Sentinel1} European Space Agency, Sentinel-1 SAR User Guide (2023).
\bibitem{GEE} N.~Gorelick \emph{et al.}, Google Earth Engine, Remote Sens.\ Environ.\ \textbf{202}, 18 (2017).
\bibitem{Ulaby2014} F.~T.~Ulaby and D.~G.~Long, Microwave Radar and Radiometric Remote Sensing (Univ.\ Michigan Press, 2014).
\end{thebibliography}

\end{document}
